\section{Pair selection can reverse the joint uncertainty comparison}
\label{exec17:s4-common-population}

The shared-control calculation becomes more informative when the action groups
have different covariate composition. Consider two equally prevalent account
strata, $X=0,1$, with the following stipulated thirty-day law:
\[
\begin{array}{c|ccc|cc}
 X&e_0&e_1&e_2&\tau_1\ ({\rm USD})&\tau_2\ ({\rm USD})\\ \hline
 0&1/2&1/4&1/4&0&12\\
 1&1/4&1/4&1/2&20&0
\end{array}
\]
Let every arm have conditional residual variance $4$ USD$^2$, with
mean-zero residuals, and retain conditional exchangeability and support.
For example, independent residuals taking values $-2,2$ with equal
probability realize that variance. Baseline means may be fixed at any
declared common values. The effects averaged over all eligible accounts are
$\theta_1=10$ and $\theta_2=6$ USD. This complete observable-law specification
permits exact score calculations; it is not a fitted bank response model.
\citep{exec17_s4_extension}

Restricting observations to a treatment/control pair changes the covariate
law. Write $B_j=1\{A\in\{0,j\}\}$ and
$s_j(x)=e_0(x)+e_j(x)$. Bayes' rule gives
\begin{equation}
 P(X=x\mid B_j=1)=\frac{P(X=x)s_j(x)}{P(B_j=1)},\qquad
 \theta_j^{\rm pair}=
       \frac{E\{s_j(X)\tau_j(X)\}}{E\{s_j(X)\}}.
 \label{exec17:s4-pair-target}
\end{equation}
For action 1, the retained mass is $5/8$ and the high-stratum share is
$2/5$, giving $\theta_1^{\rm pair}=8$ USD. For action 2, retained mass is
$3/4$ and both strata remain equally prevalent, giving
$\theta_2^{\rm pair}=6$ USD. The pair contrast for action 1 and the
full-population contrast differ even with the outcome regressions known
exactly. Dropping the other action is therefore not an innocuous way to
define the evaluation population.

This calculation explains a specific concern in the selected CausalML
implementation. Its point mean averages predictions over the full input
matrix, while its analytic uncertainty calculation uses a pair mask and
pair-row denominator. The derived $8$ is an explanatory pair-population
mean, not a claim that the package returns a point estimate of $8$; with
the true predictions and the full balanced input, its described point
average is $10$. Nor does the selected variance formula automatically
become the valid variance for either estimand under arbitrary assignment.
The full-input mean and masked ingredients require a matched statistical
argument.~\citep{exec14_causalml_primary}

For the common-population scores in Equation~\eqref{exec17:s4-shared-control},
the heterogeneity variances are $100$ and $36$, and their covariance is
$-60$ USD$^2$. The expected shared-control noise contribution is
$E(4/e_0)=12$; the active-arm contributions are $16$ and $12$.
Consequently
\begin{equation}
 \Sigma=\begin{pmatrix}128&-48\\-48&60\end{pmatrix},\qquad
 \operatorname{Var}(\phi_2-\phi_1)=284\quad {\rm USD}^2.
 \label{exec17:s4-negative-joint}
\end{equation}
Here treating the two estimates as independent gives $188$, understating
the variance of their difference. This reverses the direction of the error
in the constant-effect example, where independence gave $48$ instead of
$32$. A common control creates positive residual covariance, but opposing
heterogeneous effects can dominate it.

There is a useful separate check. The conditional difference
$\tau_2-\tau_1$ takes values $12,-20$ and has variance $256$.
In $\phi_2-\phi_1$ the control residual cancels, leaving active-arm
contributions $12+16=28$. Their sum is $284$, agreeing with the matrix
calculation. An enumeration of the two strata, three observed arms and two
residual values gives the same covariance directly from twelve observable
atoms. Thus the common-population comparison can be verified without
treating a fitted prediction matrix as a covariance model.
\citep{exec17_s4_extension}

One possible reuse design keeps the candidate package's conditional
predictions while constructing a separate common-population score evaluator.
It must retain every eligible account's covariates, original multiarm
propensities, action coding and independent sampling unit. For an externally
specified evaluation population, a declared transport or reference-sample
argument is additionally needed. Merely multiplying a pair-specific standard
error by a population-size ratio does not restore the missing heterogeneity
or joint covariance. The proposed evaluator is a development route, whose
statistical and computational behavior still needs prospective comparison.

\section{The experiment design depends on which comparison matters}
\label{exec17:s4-design-criterion}

The covariance argument also changes how a future public or prototype
experiment should be designed. Suppose a new randomized study can choose
constant probabilities $p_0,p_1,p_2>0$, summing to one. Hold the target
population fixed, take constant effects and residual variance $4$ in each
arm, and consider iid observations. If the criterion is the sum of the
two ATE influence variances, it is
\begin{equation}
 T(p)=\frac{8}{p_0}+\frac{4}{p_1}+\frac{4}{p_2}.
 \label{exec17:s4-design-trace}
\end{equation}
The control variance is charged twice because it enters both comparisons.
For a sample of size $n$, divide this criterion by $n$. It measures the
chosen precision objective, not financial utility or an experiment's total
cost.

For positive coefficients $a_j$, Cauchy--Schwarz yields
\[
 \left(\sum_j\frac{a_j}{p_j}\right)\left(\sum_jp_j\right)
       \geq\left(\sum_j\sqrt{a_j}\right)^2.
\]
Equality requires $p_j$ proportional to $\sqrt{a_j}$. Applying this with
$a=(8,4,4)$ gives
\[
 p_0=\sqrt2-1,\qquad p_1=p_2=(2-\sqrt2)/2,\qquad
 T_{\min}=24+16\sqrt2\simeq46.63\quad {\rm USD}^2.
\]
The earlier design $(1/2,1/4,1/4)$ gives $48$. Symmetry provides another
derivation: set $p_1=p_2=(1-p_0)/2$ and minimize
$8/p_0+16/(1-p_0)$. Its second derivative is
$16/p_0^3+32/(1-p_0)^3>0$, and its stationary point is $\sqrt2-1$.
Both arguments certify the global optimum for this stated criterion.
\citep{exec17_s4_extension}

If the decision is instead which active action is better, the relevant
constant-effect influence variance is
\[
 D(p)=4/p_1+4/p_2.
\]
It contains no control residual. Minimizing it alone would drive the
control probability toward zero, sacrificing precision and ultimately
support for the two comparisons against no action. If the study separately
requires $p_0\geq1/5$, this criterion is minimized at
$(p_0,p_1,p_2)=(1/5,2/5,2/5)$, giving $D=20$ but $T=60$.
The earlier design gives $D=32$ and $T=48$. Neither design dominates the
other on both precision objectives.

The distinction is useful for a bank deciding whether to fund a pilot.
An experiment intended to choose between two already acceptable servicing
actions has a different information objective from one intended to show
that either beats no contact. If effects are heterogeneous, their fixed
population covariance terms add to these formulas; if variances differ
by stratum, costs differ by arm, or randomization is clustered, the design
problem must be reformulated. The displayed optimum assumes equal per-row
resource cost and the specific independent law. It does not authorize a
real randomized bank intervention or assert power for an unspecified loss
threshold.

\section{Effect sets lead to different robust allocation objectives}
\label{exec17:s4-robust-allocation}

An effect need not be reduced to one number before allocation. Let there
be finitely many strata with known masses $w_i$, a fixed feasible allocation
set $\mathcal D$, and a rectangular net-gain class
$\mathcal G=\prod_i[\ell_i,u_i]$. This class is stipulated scientific
sensitivity information. A statistical confidence set would need a separate
simultaneous coverage argument, and correlations or common structural
restrictions could make a rectangle unnecessarily broad. For fixed known
costs, an effect interval can be translated to a net-gain interval by
subtracting the cost from both endpoints.

Write $V_g(d)=\sum_iw_i g_i d_i$, where $d_i\geq0$. The maximin policy
maximizes its worst possible value. Rectangularity and nonnegative
allocations imply
\begin{equation}
 \inf_{g\in\mathcal G}V_g(d)=\sum_iw_i\ell_i d_i.
 \label{exec17:s4-maximin}
\end{equation}
Thus maximin allocation can use the earlier resource optimizer with lower
gains. Its objective is a guarantee against the declared class, not an
estimate that every gain equals its lower endpoint.

Minimax regret asks a different question: how much would the chosen policy
lose relative to the best feasible policy if the true gain vector became
known? For a fixed $d$, its worst regret is
\begin{align}
 R(d)&=\sup_{g\in\mathcal G}\sup_{z\in\mathcal D}
                   \sum_iw_i g_i(z_i-d_i)\notag\\
 &=\sup_{z\in\mathcal D}\sum_iw_i\left\{
       u_i(z_i-d_i)_+ +\ell_i\min(z_i-d_i,0)\right\}.
 \label{exec17:s4-robust-regret}
\end{align}
The two suprema commute because they range over a product set. For each
fixed difference, the gain maximizing its contribution is the upper endpoint
for a positive difference and the lower endpoint for a negative one. The
oracle and chosen policy must share the same feasible class. Allowing the
oracle extra capacity would measure a different loss.

Take two equally prevalent strata, one resource unit per action, budget
$B=1/2$, and net gains
\[
 g_1\in[2,8],\qquad g_2\in[4,6]\quad {\rm USD}.
\]
All gains are positive, so it suffices to allocate the full budget. Write
$d=(t,1-t)$, $0\leq t\leq1$. Maximin value is
$\tfrac12\{2t+4(1-t)\}=2-t$, maximized at $t=0$ with guarantee $2$ USD.
If stratum 1 is better, regret is
$\tfrac12(1-t)(g_1-g_2)$; if stratum 2 is better, it is
$\tfrac12t(g_2-g_1)$. Therefore
\begin{equation}
 R(t)=2\max(1-t,t),\qquad t_{\rm minimax}=1/2,\quad R(1/2)=1\ {\rm USD}.
 \label{exec17:s4-robust-example}
\end{equation}
The endpoints $(g_1,g_2)=(8,4)$ and $(2,6)$ attain the two regret branches.
Their equality at $t=1/2$ and the bound
$\max(t,1-t)\geq1/2$ prove optimality. The mixed policy's guaranteed value
is $3/2$ USD. The maximin policy guarantees more, $2$ USD, while having
worst regret $2$ USD. Different objectives produce different defensible
recommendations from the same causal information.~\citep{exec17_s4_extension}

The interval midpoints both equal $5$ USD. A point-value optimizer can
therefore break that tie arbitrarily and choose either pure allocation,
each with worst regret $2$. The minimax mixture halves that worst regret.
This does not show that midpoint estimation is universally inferior:
a justified probability model for gains would support a distinct expected
utility criterion. It shows why the proposal should expose the decision
objective and preserve causal ambiguity instead of silently turning an
effect set into its midpoint. The mixture concerns population fractions;
an actual finite roster still requires an admissible implementation of
capacity and randomization.

\section{A proved finite-policy evaluation bound}
\label{exec17:s4-finite-policy}

The deterministic regret comparison can be supplemented by a transparent
statistical result in a restricted experiment. Suppose a training sample
produces a fixed finite collection of $M$ policies $d_1,\ldots,d_M$ and
outcome predictions $\widetilde m_a$. Conditional on that training sample,
evaluate the policies on $n$ new iid accounts from the declared population.
Let treatment be randomized with known $e(X)$ bounded away from zero and
one, and let costs $c(X)$ be known. Define each account's net-value score
\begin{align}
 W_j={}&d_j(X)\left\{\widetilde m_1-\widetilde m_0-c(X)
       +\frac{A}{e(X)}(Y-\widetilde m_1)
       -\frac{1-A}{1-e(X)}(Y-\widetilde m_0)\right\}.
 \label{exec17:s4-policy-score}
\end{align}
Conditional on $X$ and training, the two residual terms have expectations
$m_1-\widetilde m_1$ and $-(m_0-\widetilde m_0)$. Thus
$E(W_j\mid\text{training})=V_j=E\{d_j(X)(\tau(X)-c(X))\}$.
No correctness of the trained outcome regressions is needed for this
expectation because the assignment probability is known and correct.
This argument does not carry over unchanged to a fitted observational
propensity or an invalid scientific design.

Assume every $W_j$ lies in a known interval $[a,b]$, of width $R=b-a$.
The following elementary exponential bound gives a fully stated route to
simultaneous evaluation. For any $Z\in[0,1]$ with mean $p$, convexity of
the exponential gives
\[
 E e^{t(Z-p)}\leq (1-p)e^{-tp}+p e^{t(1-p)}.
\]
The logarithm of the right side has value and derivative zero at $t=0$;
its second derivative is the variance of a tilted Bernoulli variable,
at most $1/4$. Integrating that bound twice gives a logarithm at most
$t^2/8$. Scaling $Z=(W_j-a)/R$, independence of the evaluation accounts
therefore implies
\[
 E\exp\{\lambda\textstyle\sum_{i=1}^n(W_{ij}-V_j)\}
       \leq\exp(n\lambda^2R^2/8).
\]
The exponential Markov inequality followed by minimization at
$\lambda=4u/R^2$ yields
$P(\overline W_j-V_j\geq u)\leq\exp(-2nu^2/R^2)$.
Applying it to the negative deviation and taking a union bound over the
$M$ policies proves, conditionally on training,
\begin{equation}
 P\left(\max_j|\overline W_j-V_j|\leq
       R\sqrt{\frac{\log(2M/\alpha)}{2n}}\right)\geq1-\alpha.
 \label{exec17:s4-uniform-bound}
\end{equation}
The scores for different policies may be correlated on the same evaluation
accounts; this union argument does not require independence across policies.
It requires independence across evaluation accounts and the fixed conditional
candidate set.~\citep{exec17_s4_extension}

Let $\widehat j$ maximize evaluated mean value over that set, and let
$j^*$ maximize true value. On the simultaneous event, writing its radius
as $b_n$, one obtains
\[
 V_{j^*}-V_{\widehat j}
 \leq (V_{j^*}-\overline W_{j^*})
      +(\overline W_{\widehat j}-V_{\widehat j})\leq2b_n.
\]
An optimization tolerance $\eta$ adds $\eta$. Selection using this evaluation
set is allowed because the event holds for every member of the finite,
conditionally fixed collection. Continually inventing new policies from
these evaluation outcomes, or applying the result to an unrestricted
adaptive learner class, would exceed the proof's domain.

For an explicit monetary scale, take $e=1/2$, outcomes and trained arm
means in $[0,30]$ USD, costs in $[0,5]$ USD and $0\leq d_j\leq1$.
When $A=1$, the bracket in Equation~\eqref{exec17:s4-policy-score} is
$2Y-\widetilde m_1-\widetilde m_0-c$; when $A=0$, it is
$\widetilde m_1+\widetilde m_0-2Y-c$. Both lie in $[-65,60]$ USD;
multiplication by a policy fraction preserves that interval. Thus $R=125$.
For $M=16$, $\alpha=1/20$ and $n=10{,}000$, the simultaneous radius is
approximately $2.2468$ USD and the selected-policy regret guarantee is
$4.4936$ USD. These are proved upper bounds under the stipulated design,
not observed interval widths or measured regret.

To make the same sufficient regret guarantee at most $1$ USD, without
optimization error, requires
\begin{equation}
 n\geq 2(125)^2\log(640),\qquad n_{\rm integer}=201{,}921.
 \label{exec17:s4-evaluation-size}
\end{equation}
The conservative requirement is not a claim that this is the necessary
sample size or that a bank should purchase so many observations. A sharper
variance-sensitive analysis, fewer candidates, better score bounds or a
different acceptable loss may change the design substantially. Boundedness
is substantive: unbounded monetary outcomes, unknown truncation bias,
dependent accounts, estimated assignment probabilities or altered policy
classes need their own justified analysis. An empirical coverage study can
compare suitable procedures, but cannot be replaced by counting the exact
reference cases above.

For a product proposal, this calculation explains why public covariates and
generated interventions are useful before a costly authentic pilot. They can
exercise score conventions, evaluation splits, policy selection and complete
failure accounting while the causal truth is known by construction. The
costly study's merit should be judged against a declared precision or loss
goal. Neither the bound nor a generated study establishes the actual bank's
assignment law, feasible action set or acceptable monetary loss.

\section{What scientific information can close a discrete-score gap?}
\label{exec17:s4-gap-bounds}

The earlier two-point example shows that smoothness without a quantitative
restriction does not identify the cutoff effect. A more informative answer
can be obtained if the sponsor can defend explicit slope bounds. Suppose
the only observed scores are $-a$ and $b$, where $a,b>0$, sharp assignment
holds, and the observed means are $y_-$ and $y_+$. Let
$\mu_j(s)=E\{Y(j)\mid S=s\}$ be potential-mean extensions to the cutoff.
Stipulate that $\mu_0$ and $\mu_1$ are respectively Lipschitz with known
constants $L_0,L_1$ on the relevant intervals. Consistency then implies
\[
 \mu_0(0)\in[y_- -L_0a,y_-+L_0a],\qquad
 \mu_1(0)\in[y_+-L_1b,y_++L_1b].
\]
Subtracting the intervals gives
\begin{equation}
 \tau(0)\in[y_+-y_- -L_1b-L_0a,
             y_+-y_- +L_1b+L_0a].
 \label{exec17:s4-gap-sharp}
\end{equation}
These are sharp bounds for this two-point observed law and these separate
Lipschitz restrictions, absent further outcome or shape constraints.
For the lower endpoint choose the linear extensions
\[
 \mu_1(s)=y_++L_1(s-b),\qquad
 \mu_0(s)=y_-+L_0(s+a).
\]
They preserve the observed arm means, obey the slope bounds and attain the
lower difference at zero. Reversing both slopes attains the upper endpoint.
Convex mixtures preserve the same observations and restrictions while
attaining every intermediate difference. This is a constructive scientific
identified set, not a confidence interval produced by a smoother.
\citep{exec17_s4_extension}

For $a=b=1$, $y_-=10$, $y_+=14$ and $L_0=L_1=1$, the effect set is
$[2,6]$ USD. With action cost $3$ USD the net set is $[-1,3]$. A unit
slope restriction excludes the earlier zero-effect response system, whose
baseline slope was two; that exclusion comes from the additional scientific
restriction. At $L_0=L_1=2$, the set widens to $[0,8]$ and includes both
earlier systems. The data have not changed. The quantitative assumption
explains exactly why the conclusion changes.

More observations at the same two scores may reduce sampling error in
$y_-$ and $y_+$, but cannot remove the deterministic width
$2(L_1b+L_0a)$. Closer supported scores, justified smaller slope bounds,
intervention information in the gap or an alternative scientifically
defensible target could reduce that ambiguity. When observed means are
estimated, a valid joint confidence region for them can be propagated
through the endpoint expressions; its sampling uncertainty remains distinct
from the stipulated $L_j$. A Bayesian extrapolation can likewise be useful
with an explicit prior and model, while answering a different uncertainty
question.

The effects proposal can consequently offer concrete remedies with precise
limits: align evaluation populations, preserve joint comparisons, optimize
the intended resource and uncertainty criterion, and state which additional
information narrows an unresolved effect set. Selected software can support
parts of those routes. Their integration, runtime qualification and bank
assumptions remain explicit development and adoption obligations.
\label{exec17:s4-extension-end}
